\section{Разностная схема}

Пусть $M$ и $N$ обозначают числа интервалов по координатам $x$ и $y$.
Узлы прямоугольной сетки задаются формулами
\begin{equation}
    \begin{aligned}
        x_i & =-4+ih_x,\quad h_x=8/M,\quad 0\leq i\leq M;  \\
        y_j & =-6+jh_y,\quad h_y=12/N,\quad 0\leq j\leq N.
    \end{aligned}
\end{equation}
Полная сетка содержит $(M+1)(N+1)$ узлов.
Значения на внешней границе известны, поэтому число неизвестных равно $K=(M-1)(N-1)$.
Запись $M\times N$ обозначает числа интервалов по координатам.

\subsection{Коэффициенты около границы}

Для каждого внутреннего узла $(i,j)$ вводится ячейка
\begin{equation}
    \Pi_{ij}=[x_i-h_x/2,x_i+h_x/2]\times[y_j-h_y/2,y_j+h_y/2].
\end{equation}
Обозначим площадь пересечения этой ячейки с $D$ через $S_{ij}$.
Тогда $s_{ij}=S_{ij}/(h_xh_y)$ равно доле её площади внутри области.
Усреднение коэффициентов из~\eqref{eq:coefficients} даёт
\begin{equation}
    F_{ij}=s_{ij},\qquad Q_{ij}=\frac{1-s_{ij}}{\varepsilon^2}, \qquad \varepsilon=\max(h_x,h_y)^2.
    \label{eq:cellcoeff}
\end{equation}
В полностью внутренней ячейке $s_{ij}=1$, в полностью внешней $s_{ij}=0$.
Частично пересечённые ячейки учитывают положение границы без округления до ближайшего узла.

Площадь $S_{ij}$ вычисляется аналитическим интегрированием горизонтальных сечений.
При $y\geq0$ граница относится к кругу с центром $(0,2)$, при $y<0$ к кругу с центром $(0,-2)$.
Полуширина сечения равна $\sqrt{9-(|y|-2)^2}$ там, где подкоренное выражение неотрицательно.
Интервал интегрирования делится в точках $y=0$ и пересечения дуги с вертикальными сторонами ячейки.
На каждом участке используется первообразная
\begin{equation}
    \int\sqrt{9-t^2}\,dt
    =\frac{t\sqrt{9-t^2}+9\arcsin(t/3)}{2}.
\end{equation}
Геометрическое интегрирование не содержит дополнительного параметра дискретизации.

\subsection{Система линейных уравнений}

Вторая производная заменяется центральной разностью.
Для внутренних узлов, то есть $1\leq i\leq M-1$ и $1\leq j\leq N-1$, получается уравнение
\begin{equation}
    -\frac{v_{i+1,j}-2v_{ij}+v_{i-1,j}}{h_x^2}
    -\frac{v_{i,j+1}-2v_{ij}+v_{i,j-1}}{h_y^2}
    +Q_{ij}(v_{ij}-g_{ij})=F_{ij},
    \label{eq:stencil}
\end{equation}
где $g_{ij}=\sin(x_i+y_j)$.
Известные граничные значения переносятся в правую часть.
Матрица $\mat A\in\mathbb R^{K\times K}$ симметрична и положительно определена: к разностному оператору $-\Delta$ с условиями Дирихле добавляется неотрицательная диагональ $Q_{ij}$.
Её действие вычисляется по пятиузловому шаблону, без хранения полной матрицы.

На мелких сетках $\varepsilon^{-2}$ становится большим.
Чтобы избежать вычитания больших близких слагаемых $Q_{ij}v_{ij}$ и $Q_{ij}g_{ij}$, решается эквивалентная система для поправки $w_{ij}=v_{ij}-g_{ij}$:
\begin{equation}
    \mat A\vect w=\vect b,\qquad
    b_{ij}=F_{ij}-\lambda_h g_{ij},\qquad
    \lambda_h=\frac{4\sin^2(h_x/2)}{h_x^2}
    +\frac{4\sin^2(h_y/2)}{h_y^2}.
    \label{eq:shifted}
\end{equation}
Формула для $\lambda_h$ следует из применения центральных разностей к $\sin(x+y)$.
На внешней границе поправка равна нулю.
После решения восстанавливается $v_{ij}=w_{ij}+g_{ij}$; именно эти значения сохраняются как результат.
Жирные символы $\vect w,\vect b\in\mathbb R^K$ обозначают векторы внутренних узлов, упорядоченных по двум индексам сетки.

\section{Метод сопряжённых градиентов}

Симметрия и положительная определённость $\mat A$ позволяют использовать метод сопряжённых градиентов~\cite{cg}.
Диагональный предобуславливатель $\mat D$ уменьшает различие масштабов коэффициентов:
\begin{equation}
    D_{ij}=2/h_x^2+2/h_y^2+Q_{ij}.
\end{equation}
Здесь $D_{ij}$ обозначает диагональный коэффициент, соответствующий узлу $(i,j)$.
Вектор невязки $\vect r=\vect b-\mat A\vect w$ показывает, насколько текущее приближение удовлетворяет системе.
Вектор $\vect z=\mat D^{-1}\vect r$ получается делением каждого элемента невязки на соответствующий диагональный коэффициент.
Во всех скалярных произведениях используется сеточная мера:
\begin{equation}
    \inner{\vect a}{\vect b}=h_xh_y\sum_{i=1}^{M-1}\sum_{j=1}^{N-1}a_{ij}b_{ij},
    \qquad \normE{\vect a}=\sqrt{\inner{\vect a}{\vect a}}.
\end{equation}

Начальное приближение для поправки $\vect w^{(0)}=0$ соответствует $v^{(0)}=g$.
Вычисляются $\vect r^{(0)}=\vect b$, $\vect z^{(0)}=\mat D^{-1}\vect r^{(0)}$ и $\vect p^{(0)}=\vect z^{(0)}$.
На итерации $k=0,1,\ldots$ используются формулы
\begin{align}
    \alpha_k        & =\frac{\inner{\vect r^{(k)}}{\vect z^{(k)}}}
    {\inner{\mat A\vect p^{(k)}}{\vect p^{(k)}}},                  \\
    \vect w^{(k+1)} & =\vect w^{(k)}+\alpha_k\vect p^{(k)},\qquad
    \vect r^{(k+1)}=\vect r^{(k)}-\alpha_k\mat A\vect p^{(k)},     \\
    \vect z^{(k+1)} & =\mat D^{-1}\vect r^{(k+1)},\qquad
    \beta_k=\frac{\inner{\vect r^{(k+1)}}{\vect z^{(k+1)}}}
    {\inner{\vect r^{(k)}}{\vect z^{(k)}}},                        \\
    \vect p^{(k+1)} & =\vect z^{(k+1)}+\beta_k\vect p^{(k)}.
\end{align}
Вектор $\vect p^{(k)}$ задаёт направление изменения решения, а $\alpha_k$ величину шага вдоль него.
Итерации останавливаются при выполнении условия задания
\begin{equation}
    \normE{\vect w^{(k+1)}-\vect w^{(k)}}<\delta.
    \label{eq:stop}
\end{equation}
Для основных расчётов выбрано $\delta=10^{-8}$, а для проверки на известном решении $\delta=10^{-10}$.
Это порог изменения итерации, а не оценка ошибки исходной краевой задачи.
После остановки невязка повторно вычисляется непосредственно как $\vect b-\mat A\vect w$.

Дополнительно контролируется функционал
\begin{equation}
    J(\vect w)=\tfrac12\inner{\mat A\vect w}{\vect w}-\inner{\vect b}{\vect w} =-\tfrac12\inner{\vect b+\vect r}{\vect w}.
\end{equation}
Если его рост превышает $16\epsilon_{\rm mach}
    K\max(1,|J_k|,|J_{k+1}|)$, шаг откатывается и метод перезапускается с заново вычисленной невязкой.
Здесь $\epsilon_{\rm mach}$ обозначает машинную точность типа \code{double}; множитель $K$ учитывает накопление ошибок при суммировании.
Общее число итераций ограничено $K$, число допустимых перезапусков равно восьми.
